\chapter{Compatible formal lifts}
\label{chap:compatible-formal-lifts}

A homology class can admit a first-order lift without admitting a
second-order lift. Each later obstruction depends on the corrections
already chosen. A formal cycle therefore requires compatible choices of
primitives at every order. Four vectors exhibit a second-order failure.
Longer finite complexes place the first nonzero obstruction arbitrarily late.

A contraction chooses primitives linearly and defines successive
obstruction maps on homology. If every obstruction map vanishes, the
resulting formal series define a quasi-isomorphism. Homology concentrated
in one parity guarantees this vanishing, without bounds on degrees or
dimensions.

\section{The equation for a lift}
\label{sec:lift-equation}

Let $k$ be a field of arbitrary characteristic. We use chain degrees in
this chapter: the differential lowers degree. A graded vector space is a
family $(C_n)_{n\in\mathbb Z}$, with underlying vector space
$\bigoplus_n C_n$. The direct sum consists of finite sums of homogeneous
elements. A chain complex has linear maps
\[
 b:C_n\longrightarrow C_{n-1},\qquad b^2=0.
\]

Its cycles, boundaries, and homology are
\[
 Z_n=\ker(b:C_n\to C_{n-1}),\qquad
 \mathcal B_n=\operatorname{im}(b:C_{n+1}\to C_n),\qquad
 H_n(C,b)=Z_n/\mathcal B_n.
\]
The equation $b^2=0$ implies $\mathcal B_n\subset Z_n$.
For $c\in Z_n$, write $[c]$ for its homology class.
A primitive of $v\in C_n$ is an element $a\in C_{n+1}$ satisfying $ba=v$.

To perturb $b$ while retaining a differential, take a second map
$B:C_n\to C_{n+1}$ satisfying
\begin{equation}\label{eq:lift-mixed-identities}
 b^2=B^2=0,\qquad bB+Bb=0.
\end{equation}
These data define a mixed complex. Both maps have odd degree.
Our sign convention interchanges homogeneous factors of degrees $p,q$
with sign $(-1)^{pq}$. In particular, the anticommutator in
\eqref{eq:lift-mixed-identities} has a plus sign.

The operator $B$ induces a map
\[
 B_*:H_n(C,b)\longrightarrow H_{n+1}(C,b),\qquad [c]\longmapsto[Bc].
\]
Indeed, $bc=0$ gives $bBc=-Bbc=0$, and $B(ba)=-b(Ba)$ is a boundary.
Thus $B_*=0$ means that every $Bc$, for $c$ a cycle, has a primitive.
The definition supplies no choice of that primitive.

Introduce an even formal variable $u$ of degree $-2$.
Then $D=b+uB$ has degree $-1$, and
\[
 D^2=b^2+u(bB+Bb)+u^2B^2=0.
\]
Moving the even variable past either map contributes no sign.

To solve $Dc=0$ with arbitrarily many powers of $u$, define the space of formal series
in each total degree by
\begin{equation}\label{eq:lift-product-carrier}
 (C[[u]])_n=\prod_{j\geq0}u^j C_{n+2j}.
\end{equation}
An element is a sequence $c_j\in C_{n+2j}$, written
$\sum_{j\geq0}u^jc_j$. Equality means equality of every coefficient.
The space of polynomials uses the direct sum
$(C[u])_n=\bigoplus_{j\geq0}u^jC_{n+2j}$ and permits only finitely
many nonzero coefficients. These spaces can differ even in a single
total degree.

For $N\geq0$, let $F^N(C[[u]])_n$ consist of the series whose
coefficients with index less than $N$ vanish. These subspaces form the
neighborhoods of zero. A sequence converges in this topology when each
finite initial list of coefficients eventually becomes constant.

Every compatible collection of finite truncations determines exactly one
series by taking its coefficients. This proves completeness.
The only series in every $F^N$ is zero, which proves separation.
No analytic convergence condition is imposed.

The differential on \eqref{eq:lift-product-carrier} is
\begin{equation}\label{eq:lift-differential}
 D\left(\sum_{j\geq0}u^jc_j\right)
 =bc_0+\sum_{j\geq1}u^j(bc_j+Bc_{j-1}).
\end{equation}
The coefficient at index $j$ belongs to $C_{n+2j-1}$, as required
for total degree $n-1$. Each output coefficient uses at most two inputs.
The formula therefore defines a continuous map that preserves $F^N$.

All graded complexes below mean the family of their homogeneous parts,
or its direct sum over total degree. Multiplication by $u$ is the map
\[
 u:(C[[u]])_{n+2}\longrightarrow (C[[u]])_n,
 \qquad (c_0,c_1,\ldots)\longmapsto(0,c_0,c_1,\ldots).
\]
It is injective, has image $F^1(C[[u]])_n$, and commutes with $D$.
Thus the direct sum over total degrees is a graded $k[u]$-module.
An arbitrary scalar series in $k[[u]]$ need not act on that direct sum,
since its powers can create infinitely many total degrees.
The notation $C[[u]]$ will always mean the degreewise product
\eqref{eq:lift-product-carrier}, without a further product over degrees.

The negative cyclic homology of this mixed complex is, by definition,
\[
 HC^-_n(C)=H_n(C[[u]],D).
\]
The definition concerns the specified mixed complex and requires no
realization as the chains of an algebra. A class in $HC^-_n(C)$ lifts
$\alpha\in H_n(C,b)$ if the constant coefficient of a representative
has class $\alpha$. This defines the reduction map
\[
 \rho_n:HC^-_n(C)\longrightarrow H_n(C,b).
\]
It is well-defined: the constant coefficient of a $D$-cycle is a
$b$-cycle, and that of a $D$-boundary is a $b$-boundary.

By \eqref{eq:lift-differential}, the formal cycle equation is exactly
\begin{equation}\label{eq:lift-recurrence}
 bc_0=0,\qquad bc_j=-Bc_{j-1}\quad(j\geq1).
\end{equation}
For a prescribed cycle $c_0$, the first correction exists if and only if
$B_*[c_0]=0$. Every later correction has its own obstruction.

\begin{proposition}\label{prop:lift-successive-obstruction}
Let $N\geq0$ and suppose $c_j\in C_{n+2j}$ satisfy
\eqref{eq:lift-recurrence} through index $N$.
Then $Bc_N$ is a $b$-cycle, and the partial series extends through index
$N+1$ if and only if
\[
 [Bc_N]=0\quad\text{in }H_{n+2N+1}(C,b).
\]
\end{proposition}

\begin{proof}
If $N=0$, then $bBc_0=-Bbc_0=0$.
If $N\geq1$, the preceding equation gives
$bBc_N=-Bbc_N=B^2c_{N-1}=0$.
An extension must satisfy $bc_{N+1}=-Bc_N$, so it makes $Bc_N$ a
boundary. Conversely, if $Bc_N=ba$ with $a\in C_{n+2N+2}$,
choose $c_{N+1}=-a$. This is the required extension.
\end{proof}

For the same partial series, cancellation gives the exact identity
\[
 D\left(\sum_{j=0}^N u^jc_j\right)=u^{N+1}Bc_N.
\]
The partial series is a cycle modulo $F^{N+1}$. Its obstruction decides
whether it extends to a cycle modulo $F^{N+2}$.
The obstruction belongs to the specified partial lift.
If its last coefficient changes by a cycle $z$, its obstruction changes
by $B_*[z]$. Changes to earlier coefficients must also respect all
preceding equations. Independent finite truncations do not constitute
one compatible infinite series.

\section{A later obstruction and its torsion}
\label{sec:lift-finite-obstructions}

Let $C$ have one basis vector in each degree from zero to three:
\[
 |x|=0,\qquad |y|=1,\qquad |z|=2,\qquad |w|=3.
\]
All other components vanish. Set
\[
 bz=y,\qquad Bx=y,\qquad Bz=w,
\]
and set every unlisted value of both maps to zero.
The images of $b$ and $B$ lie in the span of $y,w$, on which both maps
vanish. Therefore $b^2=B^2=bB=Bb=0$ in every characteristic.

The $b$-cycles are spanned by $x,y,w$, and its boundaries by $y$.
Consequently $H_0(C,b)=k\cdot[x]$, $H_3(C,b)=k\cdot[w]$, and all other
homology groups vanish. Since $Bx=bz$ and $Bw=0$, the induced map $B_*$
is zero.

The class $[x]$ has no formal lift. Its constant representative is $x$,
and the equation $bc_1=-y$ forces $c_1=-z$.
The next equation would require $bc_2=w$ with $c_2\in C_4=0$.
Equivalently,
\[
 D(x-uz)=-u^2w.
\]
The first obstruction vanishes, while the second is the nonzero class
$-[w]$.

The whole completed homology records this failure. Put
\[
 x'=x-uz,\qquad y'=y+uw.
\]
Both substitutions preserve total degree. Their inverses are
$x=x'+uz$ and $y=y'-uw$. Direct computation gives
\[
 Dx'=-u^2w,\qquad Dz=y',\qquad Dy'=Dw=0.
\]

The pair generated by $z,y'$ is contractible: the degree-one map sending
$y'$ to $z$ and $z$ to zero satisfies $Dh+hD=1$ on this pair and all
its $u$-multiples. For a cycle $v$, that equation reads $v=D(hv)$,
so the pair contributes no homology.

On the remaining generators, the differential is multiplication by
$-u^2$ from the family generated by $x'$ to the family generated by
$w$. Multiplication by $u^2$ is injective, since it shifts coefficients
by two places. The homology has basis $[w],u[w]$, and $u^2[w]=0$.
Thus, as a graded $k[u]$-module,
\[
 HC^-_*(C)=\bigl(k[u]/(u^2)\bigr)[w],\qquad |w|=3.
\]

Its nonzero homogeneous parts have degree $3$ and $1$.
The quotient also equals $k[[u]]/(u^2)$, since every scalar series has a
unique remainder of degree less than two. The nonzero class $[w]$ has
torsion of exact exponent two: $u^2[w]=0$ and $u[w]\ne0$.

The first nonzero obstruction can occur arbitrarily late.
Fix $r\geq2$ and take basis vectors
\[
 x_0,\ldots,x_{r-1},\quad y_1,\ldots,y_{r-1},\quad w,
 \qquad |x_i|=2i,\quad |y_i|=2i-1,\quad |w|=2r-1.
\]
Define
\[
 \begin{aligned}
 bx_i&=y_i &&(1\leq i\leq r-1),\\
 Bx_i&=y_{i+1} &&(0\leq i\leq r-2),\\
 Bx_{r-1}&=w.
 \end{aligned}
\]
Every unlisted value is zero. Again both maps vanish on all their
images, so the mixed identities hold. Its homology is spanned by
$[x_0]$ in degree zero and $[w]$ in degree $2r-1$, and $B_*=0$.

For a degree-zero lift of $[x_0]$, each coefficient is a scalar multiple
of the single available vector $x_i$. The recurrence forces that scalar
to be $(-1)^i$. Hence
\begin{equation}\label{eq:lift-delayed-series}
 D\left(\sum_{i=0}^{r-1}(-u)^i x_i\right)
 =(-1)^{r-1}u^r w.
\end{equation}
The corrections through $u^{r-1}$ exist. The correction at $u^r$ would
require a primitive of $w$ in $C_{2r}=0$.

\begin{proposition}\label{prop:lift-delayed-torsion}
For this mixed complex,
\[
 HC^-_*(C)=\bigl(k[u]/(u^r)\bigr)[w],\qquad |w|=2r-1.
\]
Its nonzero homogeneous parts are one-dimensional in degrees
$2r-1,2r-3,\ldots,1$. The class $[w]$ has torsion of exact exponent $r$.
\end{proposition}

\begin{proof}
For $0\leq i\leq r-1$, set
\[
 x'_i=\sum_{j=i}^{r-1}(-u)^{j-i}x_j.
\]
This substitution is degree preserving, since $|u^{j-i}x_j|=2i$.
Its inverse is $x_{r-1}=x'_{r-1}$ and
$x_i=x'_i+u x'_{i+1}$ for $0\leq i<r-1$.
Cancellation of the $y$ terms gives
\[
 Dx'_0=(-1)^{r-1}u^rw,\qquad
 Dx'_i=y_i+(-1)^{r-1-i}u^{r-i}w\quad(1\leq i\leq r-1).
\]

Define $y'_i$ to be the latter right side.
It has degree $2i-1$, is a cycle, and the inverse substitution is
$y_i=y'_i-(-1)^{r-1-i}u^{r-i}w$.

The completed complex is therefore the direct sum of the finitely many
pairs $x'_i\mapsto y'_i$ for $i\geq1$ and the remaining pair
$x'_0\mapsto(-1)^{r-1}u^rw$.
Each of the first pairs is contractible by $h(y'_i)=x'_i$.
On the last pair, coefficient shifting makes the map injective.
Its image consists exactly of the multiples of $u^rw$.
The quotient has basis $[w],u[w],\ldots,u^{r-1}[w]$ in the degrees
stated. Its last basis vector is nonzero and its next multiple is zero,
which proves the assertion about the torsion exponent.
\end{proof}

For $r=2$ this is the four-vector example, with $x_1=z$ and $y_1=y$.
For $r=3$, the successive differentials are
\[
 Dx_0=uy_1,\qquad D(x_0-ux_1)=-u^2y_2,\qquad
 D(x_0-ux_1+u^2x_2)=u^3w.
\]
Given any fixed number of successive lifting equations, a larger $r$
makes all of them solvable before the nonzero obstruction appears.
Thus no fixed finite number of these equations suffices uniformly over
arbitrary mixed complexes.

\section{A contraction chooses the corrections}
\label{sec:lift-contraction}

To construct one linear lifting map for all input classes, we need
primitives that depend linearly on the input.
Over a field, complements of cycles and boundaries supply such choices.
Choose subspaces $H'_n,L_n$ with
\[
 Z_n=\mathcal B_n\oplus H'_n,\qquad
 C_n=\mathcal B_n\oplus H'_n\oplus L_n.
\]
They exist by the basis-extension and complement construction in
Lemma~\ref{lem:cc-basis-extension}.
No dimension restriction is needed.

The restriction $b:L_n\to\mathcal B_{n-1}$ is bijective.
Its kernel is $L_n\cap Z_n=0$. For surjectivity, write any preimage of a
boundary in the three summands. The first two summands have zero
differential, so the $L_n$ component is still a preimage.

Let $H_n=H_n(C,b)$ and give the graded vector space $H$ zero differential.
Define degree-zero maps $i:H\to C$ and $p:C\to H$ by including the
chosen representatives and projecting to their homology classes.
Define $h:C_n\to C_{n+1}$ as the inverse of
$b:L_{n+1}\to\mathcal B_n$ on $\mathcal B_n$, and set it to zero on
$H'_n\oplus L_n$. These maps satisfy
\begin{equation}\label{eq:lift-contraction}
 pi=1_H,\qquad bi=0,\qquad pb=0,\qquad
 bh+hb=1_C-ip,\qquad h^2=hi=ph=0.
\end{equation}

On $\mathcal B_n$ the term $bh$ is the identity and
$hb=ip=0$. On $H'_n$, the map $ip$ is the identity and $b=h=0$.
On $L_n$, the term $hb$ is the identity and $bh=ip=0$.
These checks prove the homotopy identity.
The remaining identities follow from $h(\mathcal B_n)\subset L_{n+1}$,
$h(H'_n\oplus L_n)=0$, and the definitions of $i,p$.

We call these data a contraction onto homology. More generally, over a
commutative ring $R$, a contraction onto $H=H_*(C,b)$ will mean $R$-linear
maps of these degrees satisfying \eqref{eq:lift-contraction}, normalized
by
\[
 pz=[z]\quad(z\in Z_n).
\]
This normalization makes $i$ a choice of representatives for the given
homology classes, since $pi=1$ and $bi=0$.

Such contractions need not exist over a ring. For the complex
$\mathbb Z\xrightarrow{2}\mathbb Z$ in degrees one and zero, a map
$i:\mathbb Z/2\mathbb Z\to\mathbb Z$ with $pi=1$ cannot exist:
every homomorphism into $\mathbb Z$ sends the element killed by two to
zero. Over a field, the construction above supplies the normalization.

For a cycle $z$, the contraction identity reduces to
\begin{equation}\label{eq:lift-cycle-decomposition}
 z=i(pz)+b(hz).
\end{equation}
In particular, $pz=0$ makes $hz$ a primitive of $z$.

For $\alpha\in H_n$, this suggests the corrections
\[
 a_0=i\alpha,\qquad a_j=-hBa_{j-1}\quad(j\geq1).
\]
The composite $hB$ raises degree by two, so $a_j\in C_{n+2j}$.
The resulting series is
\begin{equation}\label{eq:lift-chosen-series}
 I(\alpha)=\sum_{j\geq0}(-u)^j(hB)^ji\alpha.
\end{equation}

If the recurrence holds through $a_{j-1}$, then $Ba_{j-1}$ is a cycle.
Applying \eqref{eq:lift-cycle-decomposition} to it gives
\[
 ba_j=-bhBa_{j-1}=-Ba_{j-1}+ipBa_{j-1}.
\]
The chosen next coefficient succeeds exactly when $pBa_{j-1}=0$,
because $i$ is injective. The corresponding maps on homology are
\begin{equation}\label{eq:lift-obstruction-maps}
 pB(hB)^{j-1}i:H_n\longrightarrow H_{n+2j-1},\qquad j\geq1.
\end{equation}
The first is $B_*$, by the normalization of $p$.
The later maps require the preceding choices and have odd degree.

To allow a formal series of classes as input, put
$(H[[u]])_n=\prod_{q\geq0}u^qH_{n+2q}$ and extend $I$ by
\begin{equation}\label{eq:lift-coefficient-extension}
 I\left(\sum_{q\geq0}u^q\alpha_q\right)
 =\sum_{N\geq0}u^N
   \sum_{q=0}^N(-1)^{N-q}(hB)^{N-q}i\alpha_q.
\end{equation}
Each inner sum is finite and lies in $C_{n+2N}$.
The map is $R$-linear, preserves total degree and $F^N$, and commutes
with multiplication by $u$. Preservation of $F^N$ proves continuity.

\begin{theorem}\label{thm:lift-completed-quasi-isomorphism}
Let $(C,b,B)$ be a mixed complex over a commutative ring $R$, with a
contraction onto homology as above. The map $I$ in
\eqref{eq:lift-coefficient-extension} is a chain map from $(H[[u]],0)$
to $(C[[u]],b+uB)$ if and only if
\begin{equation}\label{eq:lift-all-obstructions-vanish}
 pB(hB)^{j-1}i=0\qquad\text{for all }j\geq1.
\end{equation}
When these conditions hold, $I$ is a continuous quasi-isomorphism:
it induces, in every degree, an isomorphism
\begin{equation}\label{eq:lift-homology-product}
 (H[[u]])_n=\prod_{j\geq0}u^jH_{n+2j}
 \xrightarrow{\ \cong\ } HC^-_n(C).
\end{equation}
\end{theorem}

\begin{proof}
Suppose first that \eqref{eq:lift-all-obstructions-vanish} holds.
For a single $\alpha\in H_n$, the initial coefficient $i\alpha$ is a
$b$-cycle. At step $j\geq1$, the hypothesis gives
\[
 pBa_{j-1}=(-1)^{j-1}pB(hB)^{j-1}i\alpha=0.
\]
Inductively, the calculation preceding the theorem therefore gives
$ba_j=-Ba_{j-1}$.
Thus $DI(\alpha)=0$. For a general input series, each coefficient in
$DI$ is a finite sum of these zero coefficients. Therefore $DI=0$.

Conversely, suppose $DI=0$. Apply this to every constant input
$\alpha\in H_n$. Its coefficient equation at index $j\geq1$ is
$ba_j+Ba_{j-1}=0$. Applying $p$ and using $pb=0$ gives
$(-1)^{j-1}pB(hB)^{j-1}i\alpha=0$.
Since $n,\alpha,j$ were arbitrary, all maps in
\eqref{eq:lift-all-obstructions-vanish} vanish.

Assume these equivalent conditions. To prove surjectivity on homology,
let $z\in(C[[u]])_n$ be a $D$-cycle and set $z^{(0)}=z$.
Suppose a residual cycle $z^{(m)}$ of total degree $n+2m$ has been
constructed. Its constant coefficient $z^{(m)}_0$ is a $b$-cycle.
Set
\[
 \alpha_m=pz^{(m)}_0\in H_{n+2m},\qquad
 q_m=hz^{(m)}_0\in C_{n+2m+1}.
\]
Equation \eqref{eq:lift-cycle-decomposition} shows that
$z^{(m)}-I(\alpha_m)-Dq_m$ has zero constant coefficient.
Here $q_m$ is regarded as a constant series. There is therefore a
unique series $z^{(m+1)}$ of degree $n+2m+2$ satisfying
\[
 z^{(m)}-I(\alpha_m)-Dq_m=u z^{(m+1)}.
\]

Its left side is a cycle. Applying $D$ and using injectivity of
multiplication by $u$ proves $Dz^{(m+1)}=0$.
This completes the induction for every $m\geq0$.

Iterating the equation gives, for each $N\geq0$,
\[
 z-I\left(\sum_{m=0}^{N-1}u^m\alpha_m\right)
   -D\left(\sum_{m=0}^{N-1}u^mq_m\right)=u^Nz^{(N)}.
\]
For $N=0$, both sums are empty and this is $z=z^{(0)}$.

The full series $\alpha=\sum_{m\geq0}u^m\alpha_m$ and
$q=\sum_{m\geq0}u^mq_m$ belong to $(H[[u]])_n$ and
$(C[[u]])_{n+1}$ respectively. For a fixed coefficient index, the right
side is zero once $N$ exceeds that index. The finite coefficient formulas
for $I$ and $D$ then give $z=I(\alpha)+Dq$.
Thus every homology class lies in the image of $I$.

For injectivity, let $\alpha=\sum_{j\geq0}u^j\alpha_j$ have total
degree $n$ and suppose $I(\alpha)=Dq$, where
$q\in(C[[u]])_{n+1}$. The constant equation is $i\alpha_0=bq_0$.
Applying $p$ gives $\alpha_0=0$, hence $bq_0=0$.

By \eqref{eq:lift-cycle-decomposition},
$q_0=i(pq_0)+b(hq_0)$. Define
\[
 q'=q-I(pq_0)-D(hq_0)\in(C[[u]])_{n+1}.
\]
Its constant coefficient is zero, and $Dq'=Dq$ because $DI=0$ and
$D^2=0$. Both $\alpha$ and $q'$ are divisible by $u$.

Write $\alpha=u\alpha'$ and $q'=u q''$, with total degrees $n+2$
and $n+3$ respectively. Since $I$ and $D$ commute with $u$, their
equation and injectivity of $u$ give $I(\alpha')=Dq''$.
The same argument applied to this equation shows $\alpha_1=0$.
Induction gives $\alpha_j=0$ for every $j\geq0$, so $\alpha=0$.
The source differential is zero, and this proves injectivity on homology.
Continuity and degree preservation were proved in
\eqref{eq:lift-coefficient-extension}.
\end{proof}

\begin{corollary}\label{cor:lift-pure-parity}
Let $k$ be a field and suppose there is $\epsilon\in\{0,1\}$ such that
$H_n(C,b)=0$ whenever $n\not\equiv\epsilon\pmod2$.
Then the quasi-isomorphism of Theorem~\ref{thm:lift-completed-quasi-isomorphism}
exists. No finite-dimensionality or boundedness hypothesis is required.
\end{corollary}

\begin{proof}
Choose the contraction using the complements already constructed.
Each map in \eqref{eq:lift-obstruction-maps} increases degree by the
odd integer $2j-1$. If its source degree has parity $\epsilon$, its
target group is zero. Otherwise its source group is zero.
Thus every obstruction map vanishes, for both choices of $\epsilon$.
The theorem applies.
\end{proof}

The coefficient construction proves this isomorphism without
interchanging homology and an inverse limit.

\section{A homotopy compatible with both differentials}
\label{sec:lift-common-contraction}

A contraction that also respects $B$ extends coefficientwise to the
perturbed complexes. Let $C$ be a mixed complex over a commutative ring
$R$, with a contraction onto $H=H_*(C,b)$.
Suppose a degree-one map $\beta:H_n\to H_{n+1}$ satisfies $\beta^2=0$
and
\begin{equation}\label{eq:lift-common-contraction}
 Bi=i\beta,\qquad pB=\beta p,\qquad Bh+hB=0.
\end{equation}
The first two equations say that inclusion and projection respect the
second operator. The last equation requires the homotopy to
anticommute with it. In particular, $\beta=B_*$: apply $p$ to $Bi\alpha$
and use the normalization on cycles.

\begin{proposition}\label{prop:lift-common-contraction}
The coefficientwise extensions of $i,p,h$ satisfy
\[
 Di=i(u\beta),\qquad pD=(u\beta)p,\qquad
 Dh+hD=1-ip.
\]
Together with $pi=1$ and $h^2=hi=ph=0$, these identities show that $i,p$
induce inverse homology isomorphisms between $(C[[u]],D)$ and
$(H[[u]],u\beta)$.
If $\beta=0$, each $i\alpha$ is a formal cycle with constant coefficients,
and $i$ induces $H[[u]]\cong HC^-_*(C)$ degree by degree.
\end{proposition}

\begin{proof}
All coefficientwise extensions preserve the filtration. Their degrees
remain zero, zero, and one. The identities $bi=0$ and $Bi=i\beta$ give
$Di=ui\beta=i(u\beta)$. Similarly, $pb=0$ and $pB=\beta p$ give
$pD=(u\beta)p$. Finally,
\[
 Dh+hD=bh+hb+u(Bh+hB)=1-ip.
\]

The identities $pi=1$ and $h^2=hi=ph=0$ hold coefficientwise as well.
Also $(u\beta)^2=u^2\beta^2=0$, so the target is a chain complex.

If $z$ is a $D$-cycle, then $pz$ is a cycle for $u\beta$, and the
homotopy identity gives $z=i(pz)+D(hz)$. This proves surjectivity of
the map induced by $i$. If $a$ is a $u\beta$-cycle and $i(a)=Dq$,
applying $p$ gives $a=u\beta(pq)$, so $a$ is a boundary in the target.
This proves injectivity. When $\beta=0$, $Di\alpha=0$ and the target
differential is zero, which proves the final assertion.
\end{proof}

The compatibility equations can hold with homology in both parities.
For example, start with any graded $R$-module $H$ and put $b=B=0$ on it.
Adjoin free generators $a,s,c,t$ with
$|a|=0$, $|s|=|c|=1$, and $|t|=2$. Set
\[
 bs=a,\qquad bt=c,\qquad Ba=c,\qquad Bs=-t,
\]
and let every other value vanish. The squares vanish.
The anticommutator on $s$ is $b(-t)+B(a)=-c+c=0$.
On $a,c,t$ its terms are both zero, proving all mixed identities.

Let $i,p$ include and project $H$, and define $ha=s$, $hc=t$, with
$h=0$ elsewhere. Then $bh+hb$ is the identity on $a,s,c,t$ and zero
on $H$. The side conditions in \eqref{eq:lift-contraction} follow from
the formulas.

On $a$,
$(Bh+hB)a=-t+t=0$. On $s,c,t,H$, both terms are zero.
Moreover, $Bi=pB=0$. Thus \eqref{eq:lift-common-contraction} holds
with $\beta=0$, although $B$ itself is nonzero.
The four extra generators contribute no completed homology.

\section{A lift that requires infinitely many coefficients}
\label{sec:lift-infinite-series}

Pure parity can remove every obstruction while forcing infinitely many
nonzero corrections. Let $C$ have basis
\[
 x_i\quad(i\geq0),\qquad y_i\quad(i\geq1),
 \qquad |x_i|=2i,\quad |y_i|=2i-1.
\]
Define $bx_0=0$, $bx_i=y_i$ for $i\geq1$, and $Bx_i=y_{i+1}$
for $i\geq0$. Both maps vanish on every $y_i$, so all mixed identities
hold. Its $b$-homology is $k\cdot[x_0]$ in degree zero and vanishes elsewhere.
Indeed, each $x_i\mapsto y_i$ for $i\geq1$ is a two-term acyclic pair.

Take $i([x_0])=x_0$, let $p$ project to $[x_0]$, and put
$hy_i=x_i$, $hx_i=0$. These maps satisfy the contraction equations.

The chosen lift is
\[
 I([x_0])=\sum_{i\geq0}(-u)^ix_i\in(C[[u]])_0.
\]
Its coefficient at $u^j$ in the differential is
$((-1)^j+(-1)^{j-1})y_j=0$ for $j\geq1$. The constant coefficient
also vanishes. Thus it is a cycle. It is not a boundary, since its
constant coefficient $x_0$ is not a $b$-boundary.
Theorem~\ref{thm:lift-completed-quasi-isomorphism} gives
$HC^-_{-2j}(C)=k\,u^j[I([x_0])]$ for $j\geq0$, with every other
homogeneous part zero.

No polynomial cycle lifts $[x_0]$. A homogeneous degree-zero polynomial
has the form $\sum_{i=0}^N\lambda_i u^ix_i$.
Its constant coefficient must have $\lambda_0=1$.
The equations for $y_1,\ldots,y_N$ force
$\lambda_i=-\lambda_{i-1}=(-1)^i$. The remaining term in its
differential is $(-1)^N u^{N+1}y_{N+1}\ne0$.
This calculation includes $N=0$, when the remaining term is $uy_1$.

There is a further difference from algebraic extension of scalars.
Let $C_{\mathrm{tot}}=\bigoplus_n C_n$ be the underlying vector space,
and regard $C_{\mathrm{tot}}[[u]]=\prod_{j\geq0}u^jC_{\mathrm{tot}}$
temporarily as an ungraded space of series. The coefficient map
\[
 C_{\mathrm{tot}}\otimes_k k[[u]]\longrightarrow C_{\mathrm{tot}}[[u]],
 \qquad v\otimes\sum_j a_ju^j\longmapsto\sum_j a_ju^jv,
\]
has image exactly those series whose coefficients lie in one
finite-dimensional subspace of $C_{\mathrm{tot}}$.

A finite sum of tensors has this property by taking the span of its
first factors. Conversely, choose a basis $v_1,\ldots,v_m$ for such a
subspace and write each coefficient uniquely as
$\sum_{\ell=1}^m a_{\ell j}v_\ell$. The series is the image of
$\sum_{\ell=1}^m v_\ell\otimes(\sum_j a_{\ell j}u^j)$.

The coefficient map is also injective. Rewrite any finite tensor sum
using a basis of the span of its first factors. If all output
coefficients vanish, linear independence makes every scalar series
in that expression zero.

The coefficients $(-1)^ix_i$ of this lift are linearly independent.
Hence the lift is outside the image of the coefficient map.
The comparison takes place in $C_{\mathrm{tot}}[[u]]$, while the lift
has total degree zero in \eqref{eq:lift-product-carrier}.

In fact, no nonzero even cycle lies in the image of this tensor product
for this example. Such an element can be written
$\sum_{i=0}^N f_i(u)x_i$. The coefficient of $y_{N+1}$ in its
differential is $u f_N(u)$. Injectivity of multiplication by $u$ gives
$f_N=0$, and descending induction gives every $f_i=0$.
Thus the choice of completion changes the existence of the required
cycle.

\section{Reduction, uniqueness, and choices}
\label{sec:lift-choices}

Reduction determines which constant class a formal class lifts.
For every mixed complex, its kernel satisfies
\begin{equation}\label{eq:lift-reduction-kernel}
 \ker\rho_n=u\,HC^-_{n+2}(C).
\end{equation}

Indeed, multiplication by $u$ produces zero constant coefficient,
so its image lies in the kernel. Conversely, take a cycle $z$ whose
constant coefficient is a boundary, say $z_0=bq_0$ with
$q_0\in C_{n+1}$. Then $z-Dq_0=u z'$ for a unique series $z'$ of
total degree $n+2$. Applying $D$ and using injectivity of $u$ on series
gives $Dz'=0$. Thus $[z]=u[z']$, proving the other inclusion.

This argument alone does not make multiplication by $u$ injective
on homology. The finite examples have $u$-torsion.

Under Theorem~\ref{thm:lift-completed-quasi-isomorphism}, the reduction
maps form short exact sequences
\begin{equation}\label{eq:lift-short-exact}
 0\longrightarrow HC^-_{n+2}(C)\xrightarrow{\ u\ }HC^-_n(C)
 \xrightarrow{\ \rho_n\ }H_n(C,b)\longrightarrow0.
\end{equation}
Surjectivity follows from $\rho_n[I(\alpha)]=\alpha$.

For injectivity at the left, write a class as $[I(\gamma)]$ using the
theorem. If its $u$-multiple is zero, then $[I(u\gamma)]=0$.
Injectivity of the isomorphism induced by $I$ gives $u\gamma=0$,
and coefficient shifting gives $\gamma=0$.
Exactness in the middle is \eqref{eq:lift-reduction-kernel}.

Let $\alpha\in\operatorname{im}\rho_n$ and choose a lift $\xi$ of $\alpha$.
The map $\eta\mapsto\xi+\eta$ is a bijection from
$u\,HC^-_{n+2}(C)$ to $\rho_n^{-1}(\alpha)$, with inverse
$\zeta\mapsto\zeta-\xi$, by \eqref{eq:lift-reduction-kernel}.
Thus the lift is unique exactly when this group is zero.

Under Theorem~\ref{thm:lift-completed-quasi-isomorphism}, uniqueness is
equivalent to $H_{n+2j}(C,b)=0$ for every $j\geq1$.
If all these groups vanish, the product formula gives a zero kernel.
Conversely, any nonzero element of one of them gives a nonzero series
with that coefficient and zero constant term. Its image through $I$
is a nonzero kernel class.

For a concrete degree calculation, suppose
\[
 H_2=k e_+,\qquad H_0=\bigoplus_{a=1}^{22}k e_a,\qquad
 H_{-2}=k e_-,
\]
and all other $H_n$ vanish. These hypotheses are realized by this
graded vector space with $b=B=0$, and the parity theorem applies to
any mixed complex with these homology groups.

In the identification supplied by a contraction, the groups have bases
\[
 \begin{array}{c|l|c}
 n&\text{basis}&\dim_k HC^-_n(C)\\\hline
 2&e_+&1\\
 0&e_1,\ldots,e_{22},\ u e_+&23\\
 -2-2q&u^q e_-,\ u^{q+1}e_1,\ldots,u^{q+1}e_{22},\ u^{q+2}e_+
       &24\quad(q\geq0)
 \end{array}
\]
All odd-degree groups and all groups above degree two vanish.
The notation for these basis vectors denotes their images under $I$.
The degree equation $|u^j\alpha_j|=|\alpha_j|-2j=n$ gives every entry.

In particular, $\rho_2$ is a canonical isomorphism, although chain
representatives of its unique inverse class can differ by boundaries.
The map $\rho_0$ has a one-dimensional kernel. Its lifts vary by a
multiple of $u[I(e_+)]$.

Even when every class lifts, changing the contraction can change the
lifted homology class. Take basis vectors
\[
 |x|=0,\qquad |y|=1,\qquad |z|=|a|=2,
 \qquad bz=y,\qquad Bx=y,
\]
with all other values zero. The mixed identities hold, and
$H_0=k\cdot[x]$, $H_2=k\cdot[a]$ are its only homology groups.

Keep $i([x])=x$ and $i([a])=a$ fixed. For $\lambda\in k$, define
\[
 \begin{gathered}
 h_\lambda y=z+\lambda a,\qquad
 h_\lambda x=h_\lambda z=h_\lambda a=0,\\
 p_\lambda x=[x],\qquad p_\lambda a=[a],\qquad
 p_\lambda z=-\lambda[a],\qquad p_\lambda y=0.
 \end{gathered}
\]

These are normalized contractions. On $x,a$, both sides of
$bh_\lambda+h_\lambda b=1-ip_\lambda$ vanish.
On $y$, both sides equal $y$. On $z$, both sides equal $z+\lambda a$.
The equations $bi=0$, $p_\lambda b=0$, and $p_\lambda i=1$ follow
on the same basis. Finally,
$p_\lambda h_\lambda y=-\lambda[a]+\lambda[a]=0$, and
$h_\lambda^2=h_\lambda i=0$.

The series from these contractions are finite:
\[
 I_\lambda([x])=x-u(z+\lambda a),\qquad I_\lambda([a])=a.
\]
Indeed, $D(x)=uy$, $D(z)=y$, $D(a)=0$, and $B(z+\lambda a)=0$.
The difference of the first lifts is
\[
 [I_\lambda([x])-I_0([x])]=-\lambda u[a].
\]
Every $D$-boundary has all coefficients in the span of $y$.
Thus $u[a]$ is nonzero, and these lifts differ whenever $\lambda\ne0$.

The underlying projections are nevertheless chain-homotopic for $b$.
Define $q:C_n\to H_{n+1}$ by $q(y)=-\lambda[a]$ and zero on the
other basis vectors. Then $p_\lambda-p_0=qb$, verified on $z$ and
zero on every other basis vector. Homotopy for the original differential
therefore does not identify the completed lift classes.

Suppose two degree-zero chain maps
\[
 I,I':(H[[u]],0)\longrightarrow(C[[u]],D)
\]
satisfy
$I-I'=DK$ for a degree-one map
$K:(H[[u]])_n\to(C[[u]])_{n+1}$. Then their values on every input differ
by a $D$-boundary. Hence they induce the same homology map.

A mixed-complex automorphism is an invertible degree-zero linear map
that commutes with both $b$ and $B$.
There is no universal choice of representatives invariant under all
such automorphisms. For an elementary example, let
\[
 C_1=k v,\qquad C_0=k x\oplus k y,\qquad bv=x,
\]
with other degrees zero and $B=0$.
Its degree-zero homology is $k\cdot[y]$.

Every linear representative map has $i([y])=y+c x$ for some $c\in k$.
For each $\lambda\in k$, define
\[
 T_\lambda(v)=v,\qquad T_\lambda(x)=x,\qquad
 T_\lambda(y)=y+\lambda x.
\]
It commutes with $b$ and $B$, has inverse $T_{-\lambda}$, and induces
the identity on homology.

Invariance of $i$ would require
$T_\lambda i([y])=i([y])$. But its left side is
$y+(c+\lambda)x$, which differs from $y+cx$ when $\lambda=1$.
Thus linear representative choices exist, while no such choice is
invariant under every mixed-complex automorphism.

\section{Multiplication and pairings impose further equations}
\label{sec:lift-additional-structures}

The constructed lift preserves linear combinations, degree, the
differential, and multiplication by $u$. A product on $C$ is additional
data. Suppose that $C$ is an associative graded algebra whose product
has degree zero, and that, for homogeneous $v,w$,
\[
 b(vw)=(bv)w+(-1)^{|v|}v(bw),\qquad
 B(vw)=(Bv)w+(-1)^{|v|}v(Bw).
\]

On homogeneous formal series, define multiplication by
\[
 \left(\sum_{i\geq0}u^iv_i\right)
 \left(\sum_{j\geq0}u^jw_j\right)
 =\sum_{N\geq0}u^N\sum_{i+j=N}v_iw_j.
\]

Every coefficient is a finite sum and has the required total degree.
The even degree of $u$ and the two displayed Leibniz rules make $D$ a
derivation of this product. Products of cycles are cycles.
If $v$ is a cycle, $v(Dq)=(-1)^{|v|}D(vq)$, and if $w$ is a cycle,
$(Dq)w=D(qw)$. Thus boundaries form a two-sided ideal among cycles,
and completed homology inherits an associative product.

Multiplicativity of a lift requires the separate equation
$I(\alpha\gamma)=I(\alpha)I(\gamma)$.

Even the contraction formula need not satisfy this equation on chains.
Let $A_0$ have basis $1,e,y$, let $A_1=kz$, and let other degrees vanish.
Make $1$ a unit and impose
\[
 e^2=e,\qquad ey=ye=y,\qquad ez=ze=z,
 \qquad y^2=yz=zy=z^2=0.
\]

These products are associative. On the span of $e,y,z$, the vector $e$
is a unit and the span of $y,z$ is a square-zero ideal.
For three factors, two ideal factors make either bracketing zero.
Otherwise removing the unit factors gives the same result.
The external unit $1$ preserves this argument.

Set $bz=y$ and let $b$ vanish elsewhere. The Leibniz rule on products
with $1,e$ follows from their unit actions. Products of two elements of
the square-zero ideal and their differentiated products are zero.
Thus $b$ is a derivation, and $B=0$ makes $A$ a mixed complex of algebras.

Its homology has basis $[1],[e]$ in degree zero, with $[e]^2=[e]$.
Choose
\[
 \begin{gathered}
 i([1])=1,\qquad i([e])=e+y,\qquad
 p(1)=[1],\quad p(e)=[e],\quad p(y)=p(z)=0,\\
 h(y)=z,\qquad h(e)=-z,\qquad h(1)=h(z)=0.
 \end{gathered}
\]

The map $bh+hb$ takes the values $0,-y,y,z$ on $1,e,y,z$ respectively,
which equal the values of $1-ip$.
Also $pi=1$, $bi=pb=0$, and $h^2=hi=ph=0$.
These maps therefore form a normalized contraction.

Since $B=0$, the lift formula is $I=i$ coefficientwise. However,
\[
 I([e]^2)=e+y,\qquad I([e])^2=e+2y.
\]
Their difference is the nonzero chain $y$ over every field.
It is a boundary, so this example concerns failure of multiplicativity
on chains.

A section of reduction can fail to preserve products even on homology.
Let $A$ have basis $1,a,c$, degrees $|1|=|a|=0$, $|c|=2$, and products
\[
 a^2=a,\qquad ac=ca=c,\qquad c^2=0,
\]
with unit $1$. Associativity follows by the same unit and square-zero
ideal argument on the span of $a,c$. Set $b=B=0$.

Define a linear section $j:A\to A[[u]]$ of reduction by
\[
 j(1)=1,\qquad j(a)=a+uc,\qquad j(c)=c.
\]
Its continuous extension $J:A[[u]]\to A[[u]]$ is
\[
 J\left(\sum_{q\geq0}u^qv_q\right)
   =\sum_{q\geq0}u^qj(v_q).
\]
Each output coefficient is a finite sum, and $J$ preserves total degree
and every $F^N$.

Its inverse sends $a$ to $a-uc$ and fixes $1,c$.
It commutes with the zero differentials and reduces to the identity
modulo $u$. Nevertheless,
\[
 J(a^2)=a+uc,\qquad J(a)^2=a+2uc.
\]
The difference $uc$ is nonzero over every field. Since the differential
is zero, it is nonzero on homology as well.

Pairing preservation is another independent equation.
On this last algebra, assume $\operatorname{char}k\ne2$ and define a
symmetric bilinear form by
$\langle a,c\rangle=\langle c,a\rangle=1$, with all other basis
pairings zero. This form vanishes unless the input degrees sum to two.

Extend it to formal series by the finite coefficient rule
\[
 \left\langle\sum_i u^iv_i,\sum_j u^jw_j\right\rangle
 =\sum_{N\geq0}u^N\sum_{i+j=N}\langle v_i,w_j\rangle
 \quad\text{in }k[[u]].
\]
Then
\[
 \langle J(a),J(a)\rangle=2u\ne0=\langle a,a\rangle.
\]

Thus the lifting properties do not imply preservation of this specified
form. The example asserts no compatibility between the form and the
algebra product. A prescribed multiplication, pairing, or other
operation requires its own compatibility equation on the lift.

\section{Formal differentials with higher corrections}
\label{sec:lift-higher-perturbations}

For a formal differential with higher corrections, each coefficient of
its square must vanish. Let $V=\bigoplus_n V^n$ be a graded $k$-vector
space, and let $\hbar$ be a formal parameter of degree zero.
Take linear maps
\[
 Q_r:V^n\longrightarrow V^{n+1}\quad(r\geq0).
\]

The space of formal series in degree $n$ is
$(V[[\hbar]])^n=\prod_{j\geq0}\hbar^j V^n$, with the topology given
by vanishing initial coefficients. Define
\[
 Q(\hbar)=\sum_{r\geq0}\hbar^rQ_r,\qquad
 Q(\hbar)\left(\sum_{j\geq0}\hbar^jv_j\right)
 =\sum_{N\geq0}\hbar^N\sum_{r=0}^N Q_r v_{N-r}.
\]
Each output coefficient is a finite sum, so this is a continuous
degree-one map. All $v_j$ in a homogeneous input have the same
cochain degree because $|\hbar|=0$.

The equation $Q(\hbar)^2=0$ holds if and only if
\begin{equation}\label{eq:lift-higher-square-zero}
 \sum_{a+b=N}Q_aQ_b=0\quad(N\geq0),
\end{equation}
where $a,b$ are nonnegative integers. If the square vanishes, applying
it to constant inputs and comparing coefficients proves these identities.
Conversely, the coefficient at $\hbar^N$ in its action on any series is
\[
 \sum_{j=0}^N\left(\sum_{a+b=N-j}Q_aQ_b\right)v_j=0,
\]
so its square vanishes on every formal series.

The first three identities are
\[
 Q_0^2=0,\qquad Q_0Q_1+Q_1Q_0=0,\qquad
 Q_1^2+Q_0Q_2+Q_2Q_0=0.
\]
The second identity makes $Q_1$ act on $H^*(V,Q_0)$: it takes cycles
to cycles and boundaries to boundaries with a minus sign.
The third says that the induced operator squares to zero.
For a $Q_0$-cycle $v$, it gives $Q_1^2v=-Q_0Q_2v$, a boundary.
It permits $Q_1^2$ to be nonzero on $V$ itself.

Four vectors realize this distinction. Let
\[
 V^0=k a,\qquad V^1=k s\oplus k c,\qquad V^2=k z,
\]
with all other degrees zero. Set
\[
 Q_0c=z,\qquad Q_1a=s,\qquad Q_1s=z,\qquad Q_2a=-c,
\]
with all other values zero, including $Q_r=0$ for $r\geq3$.

Direct evaluation gives
\[
 Q(\hbar)a=\hbar s-\hbar^2c,\qquad
 Q(\hbar)s=\hbar z,\qquad Q(\hbar)c=z,\qquad Q(\hbar)z=0.
\]
Thus $Q(\hbar)^2a=\hbar^2z-\hbar^2z=0$, and its square vanishes
on the remaining basis vectors. However,
\[
 Q_1^2a=z\ne0,\qquad (Q_0Q_2+Q_2Q_0)a=-z.
\]

The whole cohomology can be computed by putting
$v=s-\hbar c$, an invertible change of basis over $k[[\hbar]]$.
Then
\[
 Q(\hbar)a=\hbar v,\qquad Q(\hbar)v=0,\qquad
 Q(\hbar)c=z,\qquad Q(\hbar)z=0.
\]

The pair $c\mapsto z$ is contractible. On the other pair, multiplication
by $\hbar$ is injective and has cokernel $k[[\hbar]]/(\hbar)$.
Therefore
\[
 H^1(V[[\hbar]],Q(\hbar))
 =\bigl(k[[\hbar]]/(\hbar)\bigr)[v],
 \qquad H^n(V[[\hbar]],Q(\hbar))=0\quad(n\ne1).
\]

At $\hbar=0$, the cohomology has classes $[a]$ in degree zero and
$[s]$ in degree one. The class $[a]$ fails its first lifting equation,
because $Q_1a=s$ is not a $Q_0$-boundary.

For general $Q(\hbar)$, compatible corrections again have successive
obstructions, with all $Q_r$ contributing.
\begin{proposition}\label{prop:lift-general-obstruction}
Assume \eqref{eq:lift-higher-square-zero}. Let $N\geq1$ and suppose
$c_0,\ldots,c_{N-1}\in V^n$ satisfy
\[
 \sum_{r=0}^m Q_r c_{m-r}=0\quad(0\leq m<N).
\]
Then
\[
 \omega_N=\sum_{r=1}^N Q_r c_{N-r}\in V^{n+1}
\]
is a $Q_0$-cycle. The partial series extends through $\hbar^N$ if and
only if $[\omega_N]=0$ in $H^{n+1}(V,Q_0)$.
\end{proposition}

\begin{proof}
Put $c^{<N}=\sum_{j=0}^{N-1}\hbar^j c_j$.
The hypotheses make every coefficient below $\hbar^N$ in
$Q(\hbar)c^{<N}$ zero, and its coefficient at $\hbar^N$ is
$\omega_N$. The coefficient at that order in
$Q(\hbar)^2c^{<N}$ is therefore $Q_0\omega_N$:
all terms with a positive-index $Q_r$ act on a lower zero coefficient.
Since the square is zero, $Q_0\omega_N=0$.

Adding $\hbar^Nc_N$ leaves the earlier equations unchanged and changes
the coefficient at order $N$ to $\omega_N+Q_0c_N$.
If an extension exists, $\omega_N=-Q_0c_N$ is a boundary.
Conversely, if $\omega_N=Q_0a$ for $a\in V^n$, choose $c_N=-a$.
This proves both directions.
\end{proof}

The hypotheses on the preceding equations are essential even for
defining the next obstruction class. In the four-vector example,
start with $c_0=a$. Its first obstruction is $[s]\ne0$, so no $c_1$
solves the equation. If one nevertheless selects $c_1=\lambda a$,
the expression at the next order is $\lambda s-c$, whose
$Q_0$-differential is $-z\ne0$. It does not define a cohomology class.

For $Q(\hbar)=Q_0+\hbar Q_1$, the coefficient identities require
$Q_0^2=0$, $Q_0Q_1+Q_1Q_0=0$, and $Q_1^2=0$.
The same coefficient calculation applies to $b+uB$, with
$|b|=-1$, $|B|=1$, and $|u|=-2$.
